\documentclass[10pt,a4paper]{amsart}
\usepackage{amsmath,amssymb,amsthm,graphicx}
\usepackage{verbatim}
\newenvironment{reference}{\comment}{\endcomment}
%\newenvironment{reference}{}{}
\newenvironment{slogan}{\comment}{\endcomment}
\newenvironment{history}{\comment}{\endcomment}
% For commutative diagrams we use Xy-pic
\usepackage[all]{xy}

% We use 2cell for 2-commutative diagrams.
\xyoption{2cell}
% Theorem environments.
%
\theoremstyle{plain}
\newtheorem{theorem}[subsection]{Theorem}
\newtheorem{proposition}[subsection]{Proposition}
\newtheorem{lemma}[subsection]{Lemma}

\theoremstyle{definition}
\newtheorem{definition}[subsection]{Definition}
\newtheorem{example}[subsection]{Example}
\newtheorem{exercise}[subsection]{Exercise}
\newtheorem{situation}[subsection]{Situation}

\theoremstyle{remark}
\newtheorem{remark}[subsection]{Remark}
\newtheorem{remarks}[subsection]{Remarks}

\numberwithin{equation}{subsection}

% Macros
%
\def\lim{\mathop{\mathrm{lim}}\nolimits}
\def\colim{\mathop{\mathrm{colim}}\nolimits}
\def\Spec{\mathop{\mathrm{Spec}}}
\def\Hom{\mathop{\mathrm{Hom}}\nolimits}
\def\Ext{\mathop{\mathrm{Ext}}\nolimits}
\def\SheafHom{\mathop{\mathcal{H}\!\mathit{om}}\nolimits}
\def\SheafExt{\mathop{\mathcal{E}\!\mathit{xt}}\nolimits}
\def\Sch{\mathit{Sch}}
\def\Mor{\mathop{\mathrm{Mor}}\nolimits}
\def\Ob{\mathop{\mathrm{Ob}}\nolimits}
\def\Sh{\mathop{\mathit{Sh}}\nolimits}
\def\NL{\mathop{N\!L}\nolimits}
\def\CH{\mathop{\mathrm{CH}}\nolimits}
\def\proetale{{pro\text{-}\acute{e}tale}}
\def\etale{{\acute{e}tale}}
\def\QCoh{\mathit{QCoh}}
\def\Ker{\mathop{\mathrm{Ker}}}
\def\Im{\mathop{\mathrm{Im}}}
\def\Coker{\mathop{\mathrm{Coker}}}
\def\Coim{\mathop{\mathrm{Coim}}}

% Boxtimes
%
\DeclareMathSymbol{\boxtimes}{\mathbin}{AMSa}{"02}

%
% Macros for moduli stacks/spaces
%
\def\QCohstack{\mathcal{QC}\!\mathit{oh}}
\def\Cohstack{\mathcal{C}\!\mathit{oh}}
\def\Spacesstack{\mathcal{S}\!\mathit{paces}}
\def\Quotfunctor{\mathrm{Quot}}
\def\Hilbfunctor{\mathrm{Hilb}}
\def\Curvesstack{\mathcal{C}\!\mathit{urves}}
\def\Polarizedstack{\mathcal{P}\!\mathit{olarized}}
\def\Complexesstack{\mathcal{C}\!\mathit{omplexes}}
% \Pic is the operator that assigns to X its picard group, usage \Pic(X)
% \Picardstack_{X/B} denotes the Picard stack of X over B
% \Picardfunctor_{X/B} denotes the Picard functor of X over B
\def\Pic{\mathop{\mathrm{Pic}}\nolimits}
\def\Picardstack{\mathcal{P}\!\mathit{ic}}
\def\Picardfunctor{\mathrm{Pic}}
\def\Deformationcategory{\mathcal{D}\!\mathit{ef}}
\usepackage{cleveref}
\crefname{originalreference}{original reference}{original references}
\Crefname{originalreference}{Original reference}{Original references}
\makeatletter
\newlabel{algebraic-lemma-representable-transformations-property-implication}{{algebraic-lemma-representable-transformations-property-implication}{source}{}{sourcepointer-1}{}}
\newlabel{algebraic-lemma-representable-transformations-property-implication@cref}{{[originalreference][0][]algebraic-lemma-representable-transformations-property-implication}{[1][0][]source}{}{}{}}
\makeatother
\crefname{originallemma}{original lemma}{original lemmas}
\Crefname{originallemma}{Original lemma}{Original lemmas}
\makeatletter
\newlabel{lemma-first-diagonal-separated-second-diagonal-closed}{{lemma-first-diagonal-separated-second-diagonal-closed}{source}{}{sourcepointer-2}{}}
\newlabel{lemma-first-diagonal-separated-second-diagonal-closed@cref}{{[originallemma][0][]lemma-first-diagonal-separated-second-diagonal-closed}{[1][0][]source}{}{}{}}
\makeatother
\crefname{originallemma}{original lemma}{original lemmas}
\Crefname{originallemma}{Original lemma}{Original lemmas}
\makeatletter
\newlabel{lemma-properties-diagonal}{{lemma-properties-diagonal}{source}{}{sourcepointer-3}{}}
\newlabel{lemma-properties-diagonal@cref}{{[originallemma][0][]lemma-properties-diagonal}{[1][0][]source}{}{}{}}
\makeatother
\crefname{originallemma}{original lemma}{original lemmas}
\Crefname{originallemma}{Original lemma}{Original lemmas}
\makeatletter
\newlabel{lemma-properties-diagonal-representable}{{lemma-properties-diagonal-representable}{source}{}{sourcepointer-4}{}}
\newlabel{lemma-properties-diagonal-representable@cref}{{[originallemma][0][]lemma-properties-diagonal-representable}{[1][0][]source}{}{}{}}
\makeatother
\crefname{originalreference}{original reference}{original references}
\Crefname{originalreference}{Original reference}{Original references}
\makeatletter
\newlabel{spaces-more-morphisms-lemma-universally-injective-unramified}{{spaces-more-morphisms-lemma-universally-injective-unramified}{source}{}{sourcepointer-5}{}}
\newlabel{spaces-more-morphisms-lemma-universally-injective-unramified@cref}{{[originalreference][0][]spaces-more-morphisms-lemma-universally-injective-unramified}{[1][0][]source}{}{}{}}
\makeatother
\crefname{originalreference}{original reference}{original references}
\Crefname{originalreference}{Original reference}{Original references}
\makeatletter
\newlabel{spaces-morphisms-lemma-universally-closed-quasi-compact}{{spaces-morphisms-lemma-universally-closed-quasi-compact}{source}{}{sourcepointer-6}{}}
\newlabel{spaces-morphisms-lemma-universally-closed-quasi-compact@cref}{{[originalreference][0][]spaces-morphisms-lemma-universally-closed-quasi-compact}{[1][0][]source}{}{}{}}
\makeatother
\title{The separation axioms for a morphism of algebraic stacks, stated through its diagonal}
\author{The Stacks Project Authors}
\date{stacks-project:stacks-morphisms.tex:local-snapshot-20261012; GNU FDL 1.2 or later}
\begin{document}
\maketitle
\noindent\textit{Editorial note.} All delivered mathematical text is the author's own, copied from the pinned Stacks Project chapter source sources/stacks-algebraic-stacks-morphisms/pinned-original-source.tex; no statement or proof step was rewritten, reordered or removed.\par\medskip
\noindent\textit{Editorial note.} Preamble provenance: the Stacks Project keeps one single shared preamble for every chapter; this card takes it directly from the staged sources/stacks-algebraic-stacks-morphisms/preamble.tex as three verbatim spans (lines 8-12, 14-18 and 154-217), so no derived file was written and nothing was added or redefined. The wrapper class and the base packages are supplied by the assembler.\par\medskip
\noindent\textit{Editorial note.} Bibliography provenance: neither the counted statement nor its proof contains a citation (this lane rescans the whole span before writing the card), so the card carries no bibliography and no entry was reconstructed.\par\medskip
\noindent\textit{Editorial note.} Selection provenance: this card was opened under the file-based control contract (CONTROL.md), with the mechanical pre-card gate run on its statement and proof interval before the card was written.\par\medskip
\noindent\textit{Allowed standard prerequisites (editorial).}\begin{itemize}
\item The identification of the second diagonal with the identity section of the relative inertia. Used inside the counted proof (source lines 1813-1825). Allows the axioms to be phrased through the two diagonals.
\item The characterization of separations through the relative inertia. Invoked inside the counted proof (source lines 1826-1840). Lists the equivalent forms of each axiom.
\item The statements comparing the two diagonals. Used at the end of the counted proof (source lines 1841-1848). Completes the list of equivalences for each axiom.
\end{itemize}\medskip
\ifdefined\hypertarget\hypertarget{sourcepointer-1}{}\fi
\noindent\textit{Source reference.} Original reference algebraic-lemma-representable-transformations-property-implication: sources/stacks-algebraic-stacks-morphisms/pinned-original-source.tex, source lines 522--522. Navigation pointer only: the counted text refers to the Stacks item algebraic-lemma-representable-transformations-property-implication, which belongs to another chapter of the project and is not part of this card. The pointer carries the target's own label as its number because this chapter's staged snapshot does not contain the other chapter's tag table; its content is not reproduced here.\par\medskip
\ifdefined\hypertarget\hypertarget{sourcepointer-2}{}\fi
\noindent\textit{Source reference.} Original lemma lemma-first-diagonal-separated-second-diagonal-closed: sources/stacks-algebraic-stacks-morphisms/pinned-original-source.tex, source lines 1768--1768. Navigation pointer only: the counted text refers to the Stacks item lemma-first-diagonal-separated-second-diagonal-closed, which is an existing labeled result of this same chapter and is not part of this card. The pointer carries the target's own label as its number and its definition line as its span; its content is not reproduced here.\par\medskip
\ifdefined\hypertarget\hypertarget{sourcepointer-3}{}\fi
\noindent\textit{Source reference.} Original lemma lemma-properties-diagonal: sources/stacks-algebraic-stacks-morphisms/pinned-original-source.tex, source lines 344--344. Navigation pointer only: the counted text refers to the Stacks item lemma-properties-diagonal, which is an existing labeled result of this same chapter and is not part of this card. The pointer carries the target's own label as its number and its definition line as its span; its content is not reproduced here.\par\medskip
\ifdefined\hypertarget\hypertarget{sourcepointer-4}{}\fi
\noindent\textit{Source reference.} Original lemma lemma-properties-diagonal-representable: sources/stacks-algebraic-stacks-morphisms/pinned-original-source.tex, source lines 405--405. Navigation pointer only: the counted text refers to the Stacks item lemma-properties-diagonal-representable, which is an existing labeled result of this same chapter and is not part of this card. The pointer carries the target's own label as its number and its definition line as its span; its content is not reproduced here.\par\medskip
\ifdefined\hypertarget\hypertarget{sourcepointer-5}{}\fi
\noindent\textit{Source reference.} Original reference spaces-more-morphisms-lemma-universally-injective-unramified: sources/stacks-algebraic-stacks-morphisms/pinned-original-source.tex, source lines 1848--1848. Navigation pointer only: the counted text refers to the Stacks item spaces-more-morphisms-lemma-universally-injective-unramified, which belongs to another chapter of the project and is not part of this card. The pointer carries the target's own label as its number because this chapter's staged snapshot does not contain the other chapter's tag table; its content is not reproduced here.\par\medskip
\ifdefined\hypertarget\hypertarget{sourcepointer-6}{}\fi
\noindent\textit{Source reference.} Original reference spaces-morphisms-lemma-universally-closed-quasi-compact: sources/stacks-algebraic-stacks-morphisms/pinned-original-source.tex, source lines 1819--1819. Navigation pointer only: the counted text refers to the Stacks item spaces-morphisms-lemma-universally-closed-quasi-compact, which belongs to another chapter of the project and is not part of this card. The pointer carries the target's own label as its number because this chapter's staged snapshot does not contain the other chapter's tag table; its content is not reproduced here.\par\medskip
\begin{lemma}
\label{lemma-definition-separated}
Let $f : \mathcal{X} \to \mathcal{Y}$ be a morphism of algebraic stacks.
Then
\begin{enumerate}
\item $f$ is separated if and only if $\Delta_{f, 1}$ and $\Delta_{f, 2}$
are universally closed, and
\item $f$ is quasi-separated if and only if $\Delta_{f, 1}$ and $\Delta_{f, 2}$
are quasi-compact.
\item $f$ is quasi-DM if and only if $\Delta_{f, 1}$ and $\Delta_{f, 2}$
are locally quasi-finite.
\item $f$ is DM if and only if $\Delta_{f, 1}$ and $\Delta_{f, 2}$
are unramified.
\end{enumerate}
\end{lemma}

\begin{proof}
Proof of (1). Assume that $\Delta_{f, 2}$ and $\Delta_{f, 1}$ are
universally closed. Then $\Delta_{f, 1}$ is separated and universally
closed by
Lemma \ref{lemma-first-diagonal-separated-second-diagonal-closed}.
By
Morphisms of Spaces,
Lemma \ref{spaces-morphisms-lemma-universally-closed-quasi-compact}
and
Algebraic Stacks,
Lemma \ref{algebraic-lemma-representable-transformations-property-implication}
we see that $\Delta_{f, 1}$ is quasi-compact.
Hence it is quasi-compact, separated, universally closed and locally of
finite type (by
Lemma \ref{lemma-properties-diagonal})
so proper. This proves ``$\Leftarrow$'' of (1).
The proof of the implication in the other direction is omitted.

\medskip\noindent
Proof of (2). This follows immediately from
Lemma \ref{lemma-first-diagonal-separated-second-diagonal-closed}.

\medskip\noindent
Proof of (3). This follows from the fact that $\Delta_{f, 2}$ is always locally
quasi-finite by
Lemma \ref{lemma-properties-diagonal-representable}
applied to $\Delta_f = \Delta_{f, 1}$.

\medskip\noindent
Proof of (4). This follows from the fact that $\Delta_{f, 2}$ is always
unramified as
Lemma \ref{lemma-properties-diagonal-representable}
applied to $\Delta_f = \Delta_{f, 1}$ shows that
$\Delta_{f, 2}$ is locally of finite type and a monomorphism.
See
More on Morphisms of Spaces,
Lemma \ref{spaces-more-morphisms-lemma-universally-injective-unramified}.
\end{proof}

\end{document}
